% ============================================================
%  \subsection: Speech-Based Conversational Architectures
%  Parent: \section{State of the Art}
% ============================================================

\subsection{Speech-Based Conversational Architectures}
\label{subsec:speech-architectures}

The design of speech-based conversational agents revolves around a fundamental
question: how should a system map an incoming acoustic signal to a meaningful
spoken response?  Ji et al.\ provide a comprehensive taxonomy in the WavChat
survey~\cite{ji2024wavechat}, identifying three broad families of
architectures---\emph{cascaded}, \emph{half-cascade} (or semi-cascaded), and
\emph{end-to-end}---that differ in the degree to which speech understanding,
language reasoning, and speech synthesis are fused into a single model.  This
subsection reviews each family and compares their trade-offs.

% ----------------------------------------------------------
\subsubsection{Cascaded Pipeline (STT--LLM--TTS)}
\label{sssec:cascaded}

The cascaded approach decomposes the problem into three independently developed
stages: an Automatic Speech Recognition (ASR) module transcribes the user's
utterance into text, a Large Language Model (LLM) produces a textual response,
and a Text-to-Speech (TTS) engine renders that response as audio.  Because each
component communicates through a well-defined textual interface, the pipeline is
highly modular: individual stages can be upgraded, replaced, or scaled without
retraining the others.

This architecture underpins the majority of commercially deployed voice
assistants, including the speech pipelines offered by Azure, Google Cloud, and
Amazon Web Services.  In the research literature, X-Talk~\cite{lin2025xtalk}
proposes a cascaded dialogue framework targeting low-latency interaction, while
ESPnet-SDS~\cite{li2025espnetsds} provides an open-source toolkit that
implements cascaded spoken-dialogue systems with pluggable ASR, LLM, and TTS
back-ends.

The principal advantages of the cascaded pipeline are its maturity and
interpretability.  The intermediate text representation enables straightforward
logging, debugging, and content moderation---capabilities that are essential in
enterprise and safety-critical settings.  Training costs are low because each
stage can leverage pre-trained, off-the-shelf models.

The main drawback is latency.  The sequential execution of three distinct
inference steps introduces cumulative delays in the range of two to five
seconds, far exceeding the approximately 200\,ms inter-turn gap that
characterises natural human conversation~\cite{heldner2010pauses}.  Moreover,
errors compound across stages: an ASR misrecognition propagates into the LLM
prompt and may trigger an irrelevant response.  Finally, the text bottleneck
discards paralinguistic information---prosody, emotion, hesitation
markers---that is often crucial for pragmatic interpretation, and prevents
native support for full-duplex interaction.

% ----------------------------------------------------------
\subsubsection{Half-Cascade (Semi-Cascaded) Architectures}
\label{sssec:half-cascade}

Half-cascade designs seek a middle ground by fusing two of the three stages
while retaining the third as a separate component.  The WavChat
taxonomy~\cite{ji2024wavechat} distinguishes three variants, each eliminating
a different inter-stage boundary.

\emph{Speech-In fusion} (Variant~A) removes the ASR stage: the LLM ingests
audio features directly through a learnable encoder and projection layer,
preserving paralinguistic cues that text transcription would discard (e.g.,
SALMONN~\cite{tang2024salmonn}, Freeze-Omni~\cite{wang2024freezeomni}).
\emph{Speech-Out fusion} (Variant~B) absorbs the TTS stage into the LLM,
which generates discrete speech tokens alongside text, reducing output latency
(e.g., LLaMA-Omni~\cite{fang2024llamaomni}).  \emph{Inner Monologue}
(Variant~C) retains text as an internal reasoning scaffold while producing
speech output directly, attempting to combine coherent reasoning with
expressive synthesis (e.g., SpeechGPT~\cite{zhang2023speechgpt}).

These architectures reduce latency and exploit richer multimodal
representations compared to cascaded pipelines.  However, they introduce
significant trade-offs.  Aligning audio and text representation spaces
requires carefully staged training curricula, and there is a documented risk
of \emph{reasoning degradation}---the model's language abilities may decline
after fine-tuning on speech data~\cite{chen2025urobench}.  Speech quality
also tends to lag behind dedicated TTS systems, and none of these models has
been evaluated beyond dyadic interaction.

% ----------------------------------------------------------
\subsubsection{End-to-End Speech-to-Speech Models}
\label{sssec:e2e}

End-to-end architectures collapse the entire spoken-dialogue pipeline into a
single neural model that maps audio input directly to audio output.  A key
enabler of this paradigm is the family of \emph{neural audio codecs}---models
such as SoundStream, EnCodec, and Mimi---that convert continuous waveforms into
compact sequences of discrete tokens amenable to autoregressive modelling.

Moshi~\cite{defossez2024moshi} is widely regarded as the first publicly
demonstrated full-duplex speech-to-speech model, achieving a mouth-to-ear
latency of approximately 200\,ms.  Commercial systems such as GPT-4o
(OpenAI), Gemini~2.5 (Google), Qwen3-Omni (Alibaba), and Kimi Audio
(Moonshot AI) have since adopted end-to-end or near-end-to-end designs, though
their architectural details remain largely proprietary.  In the academic domain,
dGSLM extends the Generative Spoken Language Model to dialogue by modelling two
audio channels jointly.

The appeal of end-to-end models is clear: by eliminating intermediate
representations, they minimise latency and preserve paralinguistic cues---pitch
contour, speaking rate, emotional colouring---that cascaded systems discard.
They also support full-duplex interaction natively, as the model can listen and
speak simultaneously.

Nevertheless, current end-to-end systems exhibit significant limitations.
URO-Bench~\cite{chen2025urobench} demonstrates that, when evaluated on
reasoning-intensive tasks, end-to-end speech models substantially underperform
their text-only LLM counterparts, suggesting that the speech modality introduces
a reasoning bottleneck.  Training costs are high, requiring large volumes of
paired speech--speech data.  Reliability remains a concern: Lin et
al.~\cite{lin2025xtalk} observe that end-to-end speech models are ``often
unsuitable for deployment'' in production settings due to hallucination,
inconsistent output quality, and limited controllability.

% ----------------------------------------------------------
\subsubsection{Comparative Summary}
\label{sssec:arch-comparison}

Table~\ref{tab:arch-comparison} synthesises the trade-offs discussed above
across nine evaluation dimensions.

\begin{table}[ht]
  \centering
  \caption{Comparison of speech-based conversational architectures.}
  \label{tab:arch-comparison}
  \small
  \begin{tabularx}{\textwidth}{@{} l >{\centering\arraybackslash}X >{\centering\arraybackslash}X >{\centering\arraybackslash}X @{}}
    \toprule
    \textbf{Aspect} & \textbf{Cascaded} & \textbf{Half-Cascade} & \textbf{End-to-End} \\
    \midrule
    Latency
      & High (2--5\,s)
      & Medium
      & Low (${\sim}$200\,ms) \\
    Reasoning quality
      & High (text LLM)
      & Medium (risk of degradation)
      & Low~\cite{chen2025urobench} \\
    Paralinguistic info
      & Lost (text bottleneck)
      & Partially preserved
      & Fully preserved \\
    Audio quality
      & High (dedicated TTS)
      & Medium
      & Variable \\
    Full-duplex
      & Not native
      & Partial
      & Native \\
    Modularity
      & High
      & Medium
      & Low \\
    Training cost
      & Low
      & Medium
      & High \\
    Production readiness
      & High
      & Medium
      & Low~\cite{lin2025xtalk} \\
    Multi-party support
      & Demonstrated
      & Unexplored
      & Unexplored \\
    \bottomrule
  \end{tabularx}
\end{table}

% ----------------------------------------------------------
\subsubsection{The Multi-Party Gap}
\label{sssec:multiparty-gap}

A striking observation emerges from the literature: virtually all end-to-end and
half-cascade models are designed and evaluated exclusively in \emph{dyadic}
settings, i.e., one human conversing with one agent.  Multi-party spoken
interaction---where an AI agent participates in a group conversation with two or
more human speakers---remains largely unexplored in the context of neural
speech-to-speech models.  Prior work on multi-party conversational AI, such as
the studies by Addlesee~\cite{addlesee2024ari} with the ARI robot and by
Axelsson and Skantze~\cite{axelsson2025furhat} with the Furhat platform,
invariably relies on cascaded pipelines to handle the additional complexity of
speaker identification, overlapping speech, and group-level turn management.
This gap represents both a limitation of the current state of the art and an
opportunity for future research.

